
\section{The LLM--Orchestrator--Tool Boundary}
\label{sec:trust-boundary}

Tool-enabled LLM systems connect language models to conventional software
through an orchestration layer. Although frameworks differ in how they
organize execution, they share the same basic operation: the LLM produces a
structured tool call and the orchestrator translates it into software
execution. We examine this interface and identify four security weaknesses
that arise from this design.

\subsection{Generic Tool-Enabled LLM Model}
\label{sec:tool-model}

A tool-enabled system consists of an \emph{LLM}, an \emph{orchestrator}, and a
set of \emph{tools}. Each tool exposes metadata containing its name,
natural-language description, and argument schema. The orchestrator provides
this metadata to the LLM, receives a structured invocation

\[
C=\left\langle T,\{a_1=v_1,\ldots,a_n=v_n\}\right\rangle,
\]

and dispatches the arguments to the corresponding tool implementation. Tool
results are returned to the LLM, which may subsequently invoke additional
tools.

Frameworks realize this process using different control abstractions.
LangGraph, Microsoft Agent Framework, and Google ADK use graph-oriented
execution; AutoGen uses message passing; OpenAI Agents SDK and AWS Strands use
agent loops; and CrewAI combines task and event-driven workflows. Regardless
of the control abstraction, tool execution eventually reduces to the same
security-relevant path:

\[
\text{Input}
\rightarrow
\text{LLM}
\rightarrow
\text{Tool Call}
\rightarrow
\text{Orchestrator}
\rightarrow
\text{Tool}.
\]

This seemingly simple translation creates four important weaknesses.

\medskip
\noindent\textit{\underline{A1: LLMs do not sanitize tool arguments.}}
Although the LLM produces the final tool call, its arguments may originate
from user input, retrieved content, or previous tool outputs. Passing a value
through an LLM does not sanitize it. For example, a model can emit
\texttt{../../etc/passwd} for a string-valued path or preserve shell
metacharacters in a command argument. Schema validation establishes that the
call is structurally valid, not that its arguments are safe. Thus,
LLM-generated arguments remain untrusted when they reach the tool.

\medskip
\noindent\textit{\underline{A2: Similar descriptions can hijack tool selection.}}
The LLM commonly selects tools using their names and natural-language
descriptions. Consequently, an attacker can register a tool whose description
closely resembles that of a legitimate tool and cause it to compete for the
same requests. The same problem applies when an orchestrator first retrieves
semantically relevant tools before asking the LLM to select one.
ToolHijacker~\cite{toolhijacker} demonstrates such attacks against tool
retrieval and selection. Therefore, semantic similarity between tool
descriptions cannot establish tool identity or trust.

\medskip
\noindent\textit{\underline{A3: Attackers can synthesize tool calls through the LLM.}}
An attacker who cannot directly access the tool dispatcher may still influence
the LLM that produces its input. Given a tool interface, the attacker can
specify the desired arguments and induce the model to reproduce the
corresponding structured call:

\[
\text{Attacker Input}
\rightarrow
\text{LLM}
\rightarrow
\text{Tool Call}
\rightarrow
\text{Tool}.
\]

We call this \emph{Repeat-After-Me}. The invocation is syntactically produced
by the LLM, but its security-sensitive contents remain attacker controlled.
Thus, model generation does not establish that a tool call or its arguments
were independently chosen by the model.

\medskip
\noindent\textit{\underline{A4: Tool descriptions may not match their implementations.}}
The metadata consumed by the LLM and the code executed by the tool are
separate artifacts. A tool can therefore describe stronger behavior than it
actually enforces. For example, a \texttt{path} may be described as a ``file
within the workspace'' while the implementation accepts
\texttt{../../etc/passwd}. Similarly, a tool may describe a URL as remote
while allowing loopback destinations. Nothing in the orchestration layer
guarantees that the implementation satisfies the semantics advertised in its
description.

\medskip
We reproduced these weaknesses across representative orchestration frameworks
(Table~\ref{tab:framework-validation}). Despite their different control-flow
models, each ultimately relies on the same LLM-to-tool interface, exposing the
same underlying security assumptions.



Together, A1--A4 show that the LLM--orchestrator boundary provides neither
input sanitization nor assurance that the selected and executed tool behaves
as advertised. We therefore argue that \emph{LLM-generated tool arguments
must remain untrusted at the point of execution}, and that tool
implementations must be checked against the security properties implied by
their documented interfaces.

\begin{table}[t]
\centering
\small
\caption{Validation of A1--A4 across representative LLM orchestration
frameworks. A checkmark denotes a reproduced weakness.}
\label{tab:framework-validation}
\begin{tabular}{llcccc}
\toprule
\textbf{Framework} & \textbf{Execution Model} &
\textbf{A1} & \textbf{A2} & \textbf{A3} & \textbf{A4} \\
\midrule
LangGraph~\cite{langgraph}
    & Graph-based exec.
    & $\checkmark$ & $\checkmark$ & $\checkmark$ & $\checkmark$ \\
MAF~\cite{microsoft-agent-framework}
    & Graph-based exec.
    & $\checkmark$ & $\checkmark$ & $\checkmark$ & $\checkmark$ \\
ADK 2.0~\cite{google-adk2}
    & Graph-based exec.
    & $\checkmark$ & $\checkmark$ & $\checkmark$ & $\checkmark$ \\
AutoGen~\cite{autogen}
    & Actor Model
    & $\checkmark$ & $\checkmark$ & $\checkmark$ & $\checkmark$ \\
Strands~\cite{strands}
    & Model--tool agent loop
    & $\checkmark$ & $\checkmark$ & $\checkmark$ & $\checkmark$ \\
CrewAI~\cite{crewai}
    & Task/Event Loop
    & $\checkmark$ & $\checkmark$ & $\checkmark$ & $\checkmark$ \\
\bottomrule
\end{tabular}
\end{table}